\chapter{Phase 3 --- Verification}\label{ch:verification}

\weeks{7--12, alongside Phase 2} \quad Milestone M2

In industry, verification typically takes more effort than design. That is not
inefficiency. A bug found in simulation costs minutes. Found on the FPGA, it
costs days. Found in silicon, it can cost a re-spin. This project has an
unusual advantage that makes verification both rigorous and pleasant: a
\emph{bit-accurate golden model} written in Python. The question is never
\enquote{does this look right?} but \enquote{are the integers identical?}

\begin{whybox}
Bit-exactness against the golden model is what makes your hardware results
meaningful. Once the RTL is bit-exact, every quality number from Phase~1 (RQ1,
RQ2) is also a statement about the hardware. And verification skills (cocotb,
coverage, assertions, formal) are among the most employable things this project
teaches.
\end{whybox}

\begin{center}\small
\begin{tabularx}{\textwidth}{@{}LL@{}}
\toprule
\textbf{Inputs} & \textbf{Outputs} \\
\midrule
Golden model and vectors (\code{spec-v1}); RTL blocks as they are written
(\cref{ch:rtl}); Verilator, Icarus and cocotb from Phase~0
& A verification plan; one testbench per block plus a top-level test; a
regression in CI; a coverage summary; formal proofs for small blocks \\
\bottomrule
\end{tabularx}
\end{center}

% =========================================================================
\section{The verification plan}\label{sec:verif:plan}

Write the plan before the tests, in \file{docs/verification_plan.md}. It lists
every feature that must work, how it will be tested, and what counts as passing
(\cref{tab:verif:plan}). Review it at every milestone. A feature with no test is
a feature you are hoping works.

\begin{table}[htbp]
\centering\footnotesize
\caption{An excerpt of a verification plan. Extend it to every block and feature.}
\label{tab:verif:plan}
\begin{tabularx}{\textwidth}{@{}lLLL@{}}
\toprule
ID & Feature & Test & Pass criterion \\
\midrule
V1 & \code{qmul\_sat} rounding and saturation & Exhaustive (narrow widths) or $10^6$ random vectors & 0 mismatches against \code{fixedpoint.mul} \\
V2 & \code{recip\_unit} & Full input range or dense random & 0 mismatches against the golden reciprocal \\
V3 & \code{min\_tracker} across sub-window ends & 200 frames of golden vectors & $S_\text{min}$ bit-exact every frame \\
V4 & 5-tap smoothing at spectrum edges & Impulse at bins 0, 1, $K-2$, $K-1$ & Bit-exact with \code{np.convolve(mode="same")} semantics \\
V5 & \code{fft\_r2\_mem} with block floating point & Impulse, tones, full-scale noise & Bit-exact data and exponent \\
V6 & Parameter commit at a frame boundary & Commit in the middle of a frame & Old parameters until the boundary, new after; no mixed frame \\
V7 & Pre-roll & First \reg{PREROLL\_FRAMES} frames & Pass-through, then seeded state equals golden \\
V8 & End to end & Synthetic set + VoiceBank utterances & Bit-exact output; identical PESQ/STOI \\
V9 & Deadline & Worst-case frame & \reg{CYCLES\_LAST} below budget; no \code{FRAME\_OVERRUN} \\
\bottomrule
\end{tabularx}
\end{table}

% =========================================================================
\section{How the testbenches are built}\label{sec:verif:arch}

Every testbench has the same four parts, which mirror industrial methodologies
such as UVM in a lightweight Python form:
\begin{itemize}
  \item a \textbf{driver}, which turns test data into pin activity (valid/ready
        handshakes, samples, bus writes);
  \item a \textbf{monitor}, which watches outputs and turns them back into data;
  \item a \textbf{reference model}, which is your golden model, imported
        directly;
  \item a \textbf{scoreboard}, which compares expected and actual results and
        reports the \emph{first} mismatch with enough context to debug it.
\end{itemize}
Because cocotb testbenches are Python~\cite{cocotb}, the reference model is
not a re-implementation. It is the same code that produced your Phase~1 results.

% =========================================================================
\section{Steps}\label{sec:verif:steps}

\begin{stepbox}[label=st:verif:first]{Your first real testbench: \texttt{recip\_unit}}
The listing below is a complete pattern you will reuse for every pipelined
block. Note the separate driver and monitor coroutines. The monitor samples in
the \code{ReadOnly} phase after the clock edge, when all signals have settled.

\begin{lstlisting}[style=py]
import numpy as np
import cocotb
from cocotb.clock import Clock
from cocotb.triggers import ClockCycles, ReadOnly, RisingEdge, with_timeout
from golden.recip import recip_model          # the same code golden/ uses

async def drive(dut, values):
    for v in values:
        await RisingEdge(dut.clk)
        dut.in_data.value = int(v)
        dut.in_valid.value = 1
    await RisingEdge(dut.clk)
    dut.in_valid.value = 0

async def monitor(dut, out, n):
    while len(out) < n:
        await RisingEdge(dut.clk)
        await ReadOnly()
        if dut.out_valid.value == 1:
            out.append(int(dut.out_data.value))

@cocotb.test()
async def recip_matches_golden(dut):
    cocotb.start_soon(Clock(dut.clk, 10, "ns").start())
    dut.rst.value, dut.in_valid.value = 1, 0
    await ClockCycles(dut.clk, 3)
    dut.rst.value = 0
    rng = np.random.default_rng(1234)
    stim = [1, 2, 3, 2**15, 2**16 - 1] + list(rng.integers(1, 2**16, 5000))
    got = []
    mon = cocotb.start_soon(monitor(dut, got, len(stim)))
    await drive(dut, stim)
    await with_timeout(mon, 100, "us")
    bad = [(v, recip_model(v), g) for v, g in zip(stim, got) if recip_model(v) != g]
    assert not bad, f"{len(bad)} mismatches; first (in, expected, got): {bad[:5]}"
\end{lstlisting}

\textbf{Check:} the test passes on Verilator and on Icarus. Then deliberately
break the RTL (flip a rounding constant) and confirm the test fails with a
useful message. A test that cannot fail is not a test. \textbf{Effort:} 1 day
(with the RTL).
\end{stepbox}

\begin{stepbox}[label=st:verif:vectors]{Golden-vector plumbing}
Write \file{tb/common/vectors.py} to load the files exported by
\file{golden/vectors.py}: per test signal, per probe point, per frame, as
integer arrays. Add a scoreboard helper that compares an RTL probe stream with
its golden array. On the first mismatch it reports the frame, the bin, the
expected and actual values, and the five previous values.

\textbf{Check:} the helper reports a deliberately injected mismatch at the right
frame and bin. \textbf{Effort:} half a day.
\end{stepbox}

\begin{stepbox}[label=st:verif:binpipe]{A streaming testbench for \texttt{bin\_pipeline}}
Drive one bin per cycle from the golden spectra and compare the outputs, and
the internal probes exposed through debug ports, frame by frame. Add random
\emph{back-pressure}: hold \sig{out_ready} low on random cycles. That checks the
pipeline stalls without losing or duplicating a bin, if your design has
back-pressure at all. If it is purely schedule-driven, test the schedule
instead.

\textbf{Check:} bit-exact on all synthetic signals and at least one utterance,
including a parameter commit. \textbf{Effort:} 2--3 days.
\end{stepbox}

\begin{stepbox}[label=st:verif:top]{Frame-level test of \texttt{ise\_top} with real audio}
Feed a full VoiceBank utterance through \code{ise\_top} and collect the output
samples. Compare them bit for bit with the golden model. Then compute PESQ, STOI
and SI-SDR on the RTL output. They must equal the golden model's numbers
exactly, and that equality is the evidence connecting Phase~1 to the hardware.

\textbf{Check:} bit-exact output; identical metrics. Note the simulation time
per second of audio for your report. \textbf{Effort:} 2 days.
\end{stepbox}

\begin{tipbox}[Keep Python out of the per-cycle loop]
A full utterance is millions of clock cycles. A cocotb coroutine that wakes on
every clock edge makes the simulation slow. Wake on \emph{events} instead: a
sample strobe, a \sig{frame_done} pulse, or a FIFO becoming non-empty. Let the
simulator run the cycles in between at full speed.
\end{tipbox}

\begin{stepbox}[label=st:verif:corners]{Directed corner cases}
Real speech rarely exercises extremes. Write one directed test per item:
\begin{itemize}
  \item silence: all zeros, which tests division by tiny noise estimates and
        the $\varepsilon$ floors;
  \item a full-scale sine, which tests saturation in the FFT and the power
        computation;
  \item a clipped square wave and DC, which test bin 0 and the block exponent;
  \item an impulse and a chirp;
  \item a parameter commit in mid-frame (V6), two commits in quick succession,
        and a commit during pre-roll;
  \item a soft reset in the middle of the stream, then re-seeding;
  \item a bus read and write during frame processing.
\end{itemize}

\textbf{Check:} every item is in the verification plan with a passing test.
\textbf{Effort:} 2--3 days.
\end{stepbox}

\begin{stepbox}[label=st:verif:random]{Constrained-random tests}
Generate random but \emph{legal} stimuli: noise at random levels and spectral
shapes, random parameter sequences within the register ranges, and commits at
random times. Seed every run and print the seed, so any failure can be
reproduced exactly.

\textbf{Check:} 100 random seeds pass. Any failure becomes a new directed test
before it is fixed. \textbf{Effort:} 1--2 days.
\end{stepbox}

\begin{stepbox}[label=st:verif:coverage]{Measure coverage}
\begin{itemize}
  \item \textbf{Code coverage.} Build with Verilator's \code{--coverage} option
  and summarise with \code{verilator\_coverage}~\cite{ch12-13:verilatorguide}.
  Look for lines and toggles never exercised.
  \item \textbf{Functional coverage.} Count events that matter: every gain-table
  index used; saturation in each block; each number of block-floating-point
  shifts; commits at and away from frame boundaries; tracker sub-window ends.
  Use simple counters in Python, or the \code{cocotb-coverage} package.
\end{itemize}

\textbf{Check:} every un-hit item is either covered by a new test or explained
in the plan (for example, as unreachable by design). \textbf{Effort:} 1--2 days.
\end{stepbox}

\begin{stepbox}[label=st:verif:formal]{Assertions and formal proofs for small blocks}
Add immediate assertions inside the RTL for invariants. Verilator checks them
in simulation when built with assertions enabled:

\begin{lstlisting}[style=sv]
`ifdef FORMAL_OR_SIM_CHECKS
  always_comb begin
    assert (s_min <= running_d);
    for (int i = 0; i < U; i++) assert (s_min <= store_d[i]);
  end
`endif
\end{lstlisting}

For small blocks such as \code{min\_tracker}, the AXI-Lite register file and any
FIFO, go further: \emph{prove} the properties for all inputs with SymbiYosys,
the formal front end in the OSS CAD Suite~\cite{ch12-13:sbydocs}:

\begin{lstlisting}[style=plain,caption={\texttt{tb/formal/min\_tracker.sby} (adapt to the current SymbiYosys documentation).}]
[options]
mode prove

[engines]
smtbmc

[script]
read -formal -D FORMAL_OR_SIM_CHECKS min_tracker.sv
prep -top min_tracker

[files]
../../rtl/min_tracker.sv
\end{lstlisting}

\textbf{Check:} the proof passes. A broken variant, such as one using the wrong
comparison, produces a counterexample trace. \textbf{Effort:} 2 days, and it is
worth it: formal skills stand out on a CV.
\end{stepbox}

\begin{stepbox}[label=st:verif:ci]{Regression in CI}
Add a CI job that installs Verilator (the OSS CAD Suite is convenient for this),
runs every cocotb test and the formal proofs, and archives the coverage summary.
Keep the full-utterance tests in a nightly or manual job if they are slow.

\textbf{Check:} every push runs the fast regression. The badge is green at M2.
\textbf{Effort:} half a day.
\end{stepbox}

% =========================================================================
\section{When the RTL and the model disagree}\label{sec:verif:debug}

Mismatches are where you learn the most. Debug them systematically:
\begin{enumerate}
  \item \textbf{Find the first divergence.} Compare the internal probes block
        by block, in pipeline order. The first probe that differs, in the
        earliest frame and lowest bin, is where the bug is, or just upstream
        of it.
  \item \textbf{Shrink the case.} Reproduce it with the shortest input that
        fails, often a single frame of a synthetic signal.
  \item \textbf{Look at waves.} Dump an FST trace for just that window, with
        Verilator's tracing options, and inspect it in GTKWave or Surfer next to
        the golden values.
  \item \textbf{Decide who is wrong.} Sometimes it is the model. A float slipped
        into a fixed-point path, or the model rounds differently from the
        specification. Fix the specification first, then whichever side
        disagrees with it.
  \item \textbf{Add a regression test} for the bug before fixing it.
\end{enumerate}

\begin{pitfallbox}[Pitfalls in Phase 3]
\begin{itemize}
  \item \textbf{Comparing with tolerances.} \enquote{Within 1 LSB} hides
        systematic rounding bugs that add up over frames. Require exact
        equality.
  \item \textbf{Testing only with speech.} Speech never saturates the FFT or
        produces exact half-LSB ties. Synthetic and random signals do.
  \item \textbf{Reading outputs at the wrong time.} Sample after the edge, in
        \code{ReadOnly}. Drive on the edge, never in \code{ReadOnly}.
  \item \textbf{Unseeded randomness.} An irreproducible failure is almost
        worthless.
  \item \textbf{Verification as an afterthought.} Each block's test is written
        together with the block. A week of untested RTL is a week of hidden
        bugs.
\end{itemize}
\end{pitfallbox}

\begin{readbox}
cocotb documentation and examples~\cite{ch12-13:cocotbdocs}. Verilator user's
guide, for lint, tracing and coverage~\cite{ch12-13:verilatorguide}.
SymbiYosys~\cite{ch12-13:sbydocs}. For verification methodology in
SystemVerilog, useful even though you test in Python:
Spear and Tumbush~\cite{ch12-13:spear2012svverif}. Formal verification of real
designs, step by step: the ZipCPU blog~\cite{ch12-13:zipcpu}.
\end{readbox}

\begin{donebox}[Milestone M2 is reached (with Phase 2) when]
\begin{checklist}
  \item The verification plan covers every block and feature, and every entry has a passing test.
  \item \code{ise\_top} is bit-exact with the golden model on all golden vectors and full utterances, with identical metrics.
  \item Directed corner cases and 100 constrained-random seeds pass.
  \item Coverage is measured and every gap is closed or explained.
  \item Formal proofs pass for at least \code{min\_tracker} and the register file.
  \item The regression runs in CI and is green.
\end{checklist}
\end{donebox}
